\subsection{平坦模}

定义 1. —— 称 A-模 E 为平坦的，如果对于右 A-模与同态的每个正合序列，

\[
\mathbf {M} ^ {\prime} \xrightarrow {u} \mathbf {M} \xrightarrow {v} \mathbf {M} ^ {\prime \prime},\tag{11}
\]

由Z-线性映射构成的序列

\[
\mathbf {M} ^ {\prime} \otimes_ {\mathbf {A}} \mathrm{E} \xrightarrow {u \otimes 1} \mathbf {M} \otimes_ {\mathbf {A}} \mathrm{E} \xrightarrow {v \otimes 1} \mathbf {M} ^ {\prime \prime} \otimes_ {\mathbf {A}} \mathrm{E}\tag{12}
\]

是正合的。

命题 3. —— A-模 E 为平坦的，当且仅当对于每个右 A-模的单射 A-同态 \(u: \mathbf{M}' \to \mathbf{M}\) ，映射 \(u \otimes 1: \mathbf{M}' \otimes_{\mathbf{A}} \mathbf{E} \to \mathbf{M} \otimes_{\mathbf{A}} \mathbf{E}\) 是单射。

如果 E 是平坦的且 \(u: \mathbf{M}' \to \mathbf{M}\) 是单射，序列 \(0 \longrightarrow \mathbf{M}' \stackrel{u}{\longrightarrow} \mathbf{M}\) 是正合的，因此也有

序列 0 \(\longrightarrow\) M' \(\otimes_{\mathbf{A}}\) E \(\xrightarrow{u\otimes1}\) M \(\otimes_{\mathbf{A}}\) E，且 \(u\otimes1\) 是单射。反之，考虑正合序列（11）；设 M'\,' \(_{1}\) = v(M)，并令 i : M'\,' \(_{1}\) → M'\,' 为典范单射

和 \(p:\mathbf{M}\to\mathbf{M}_{1}^{\prime\prime}\) 为映射 \(m\mapsto v(m)\) 。序列 \(\mathbf{M}^{\prime}\xrightarrow{u}\mathbf{M}\xrightarrow{p}\mathbf{M}_{1}^{\prime\prime}\longrightarrow0\)

是正合的；由 II，第～58 页，命题 5，序列 \(\mathbf{M}' \otimes_{\mathbf{A}} \mathbf{E} \xrightarrow{u \otimes 1} \mathbf{M} \otimes_{\mathbf{A}} \mathbf{E} \xrightarrow{p \otimes 1} \mathbf{M}_{1}'' \otimes_{\mathbf{A}} \mathbf{E}\) 因此是正合的。此外，我们有 \(v = i \circ p\) ，因此 \((v \otimes 1) = (i \otimes 1) \circ (p \otimes 1)\) ；若 E 满足命题中的条件，则 \(i \otimes 1\) 是单射，因此

\[
\operatorname{Ker} (v \otimes 1) = \operatorname{Ker} (p \otimes 1) = \operatorname{Im} (u \otimes 1)
\]

且序列 (12) 是正合的。

命题4. ---（i）设 \((\mathbf{E}_i)_{i\in \mathbf{I}}\) 是 A-模的一个族， \(\mathrm{E} = \bigoplus_{i\in \mathbf{I}}\mathrm{E}_i\) 它们的直和

为使 A-模 E 是平坦的，其充要条件是每一个 \(E_{i}\) 都是平坦的。

\begin{enumerate}
\def\labelenumi{(\roman{enumi})}
\setcounter{enumi}{1}
\tightlist
\item
  设 I 是一个右有向预序集， \((\mathrm{E}_{\alpha}, f_{\beta\alpha})\) 是关于 I 的 A-模归纳系， \(E = \varinjlim E_{\alpha}\) 为其归纳极限。若每个 A-模 \(E_{\alpha}\) 是平坦的，则 E 是平坦的。
\end{enumerate}

设 \(M^{\prime} \rightarrow M \rightarrow M^{\prime\prime}\) 是右A-模的一个正合序列。

\begin{enumerate}
\def\labelenumi{(\roman{enumi})}
\item
  为使序列 \(\bigoplus_{i\in I}(\mathbf{M}^{\prime}\otimes_{\mathbf{A}}\mathbf{E}_{i})\to \bigoplus_{i\in I}(\mathbf{M}\otimes_{\mathbf{A}}\mathbf{E}_{i})\to \bigoplus_{i\in I}(\mathbf{M}^{\prime \prime}\otimes_{\mathbf{A}}\mathbf{E}_{i})\) 正合，必要且充分的是每个序列 \(\mathbf{M}^{\prime}\otimes_{\mathbf{A}}\mathbf{E}_{i}\to \mathbf{M}\otimes_{\mathbf{A}}\mathbf{E}_{i}\to \mathbf{M}^{\prime \prime}\otimes_{\mathbf{A}}\mathbf{E}_{i}\) 都正合（II，第 13 页，命题 7），这就证明了 (i)，因为 \(\oplus (\mathbf{M}\otimes_{\mathbf{A}}\mathbf{E}_{i})\) 典范地等同于 \(\mathbf{M}\otimes_{\mathbf{A}}\mathbf{E}\) （II，第 61 页，命题 7）。
\item
  根据假设，每个序列 \(\mathbf{M}' \otimes_{\mathbf{A}} \mathbf{E}_i \to \mathbf{M} \otimes_{\mathbf{A}} \mathbf{E}_i \to \dot{\mathbf{M}}'' \otimes_{\mathbf{A}} \mathbf{E}_i\) 都是正合的，因此序列 \(\mathbf{M}' \otimes_{\mathbf{A}} \mathbf{E} \to \mathbf{M} \otimes_{\mathbf{A}} \mathbf{E} \to \mathbf{M}'' \otimes_{\mathbf{A}} \mathbf{E}\) 也是正合的，因为取归纳极限与张量积可交换（II，第 93 页，命题 7）并保持正合性（II，第 91 页，命题 3）。
\end{enumerate}

例。--- 1) 显然 \(A_{s}\) 是一个平坦A-模；于是由命题4(i)可知，每个自由A-模，更一般地每个投射A-模，都是平坦的（另见II，第63页，推论6）。

* 反之，每个有限表示的平坦 A-模都是投射的（第 5 节）。 *

\begin{enumerate}
\def\labelenumi{\arabic{enumi})}
\setcounter{enumi}{1}
\item
  由命题 4 (ii)，每个作为自由 A-模的有向归纳系的归纳极限的 A-模都是平坦的。我们将在第 6 节证明一个逆命题。
\item
  若A是半单的，则每个A-模都是投射的（VIII，§5，第1节，命题1），从而是平坦的。
\item
  * 若A是阿廷局部环（不一定交换），则一个A-模是平坦的当且仅当它是自由的（AC II，§3，第2节，命题5的推论2）。 *
\item
  若A是整环，则A的分式域K是平坦A-模（II，第118页，命题27）。
\item
  * 在AC II和III中，我们将研究当A为交换环时平坦A-模的两个重要例子：分式环 \(S^{-1}\) A，以及当A为诺特环时，A关于J-adic拓扑的分离完备化。 *
\item
  设 \(a \in \mathbf{A}\) 使得映射 \(a_{\mathbf{A}}: x \mapsto ax\) 从 \(\mathbf{A}\) 到 \(\mathbf{A}\) 是单射（“ \(a\) 不是左零因子”）。如果 \(\mathbf{E}\) 是平坦 \(\mathbf{A}\) -模，则位似 \(a_{\mathbf{E}}\) 是单射，因为它等同于 \(a_{\mathbf{A}} \otimes 1: \mathbf{A}_{d} \otimes_{\mathbf{A}} \mathbf{E} \to \mathbf{A}_{d} \otimes_{\mathbf{A}} \mathbf{E}\) 。特别地，如果 \(\mathbf{A}\) 是整环，则每个平坦 \(\mathbf{A}\) -模都是无挠的。反之，如果 \(\mathbf{A}\) 是主理想环，则每个无挠 \(\mathbf{A}\) -模都是平坦的：事实上，如果 \(\mathbf{A}\) -模 \(\mathbf{E}\) 是无挠的，则 \(\mathbf{E}\) 的每个有限生成子模都是自由的（第七章，第4节，第4号，定理4的推论2），并且 \(\mathbf{E}\) 是平坦子模的有向递增并，因此是平坦的（命题4（ii））。
\item
  设 B 为一个环，且 \(\rho: \mathbf{A} \to \mathbf{B}\) 一个同态。若 E 是平坦的 A-模，则 B-模 \(\mathbf{E}_{(\mathbf{B})} = \mathbf{B} \otimes_{\mathbf{A}} \mathbf{E}\) 是平坦的。事实上，设 \(u: \mathbf{N}' \to \mathbf{N}\) 是右 B-模的一个单同态；则 \(u \otimes_{\mathbf{B}} 1_{\mathbf{E}_{(\mathbf{B})}}\) 典范地等同于同态 \(u \otimes_{\mathbf{A}} 1_{\mathbf{E}}: \mathbf{N}' \otimes_{\mathbf{A}} \mathbf{E} \to \mathbf{N} \otimes_{\mathbf{A}} \mathbf{E}\) ，若 E 是平坦的，则该同态是单的。
\item
  假设 \(\mathbf{A} = \mathbf{K}[\mathbf{X}, \mathbf{Y}]\) ，其中 \(\mathbf{K}\) 是一个域。则极大理想 \(m\) 由 \(\mathbf{X}\) 和 \(\mathbf{Y}\) 生成，它是一个无挠 A-模，但不是平坦的。事实上，考虑环 \(\mathbf{B} = \mathbf{A}/(\mathbf{Y})\) ，它同构于 \(\mathbf{K}[\mathbf{X}]\) ，因此是一个整环。B-模 \(m_{(B)}\) 同构于 \(m/Ym = (X, Y)/(XY, Y^2)\) ，其中 \(Y\) 的类是挠元。因此， \(m_{(B)}\) 不是平坦的B-模，从而 \(m\) 不是平坦的。
\item
  假设 A 是交换的。设 B 为代数 \(\mathbf{A}[\mathbf{X}_1, \ldots, \mathbf{X}_n] / (\mathbf{P})\) ，其中 P 是非零多项式。对于 A 的每个素理想 p，记 \(\kappa(p)\) 为整环 \(\mathbf{A} / p\) 的分式域， \(\mathbf{E}(p)\) 为代数 \(\kappa(p) [\mathbf{X}_1, \ldots, \mathbf{X}_n]\) ，并记 \(\mathbf{P}(p)\) 为 P 在典范映射下在 \(\mathbf{E}(p)\) 中的像。
\end{enumerate}

可以证明，要使 B 成为平坦的 A-模，只需 \(P(p) \neq 0\) 对 A 的每个素理想 p 成立。若 A 是整环，则此条件也是必要的。

* 用几何语言来说，考虑投影 \(\pi: \text{Spec}(B) \to \text{Spec}(A)\) 。对每个 \(p \in \text{Spec}(A)\) ，纤维 \(\pi^{-1}(p)\) 等同于子簇 \(V_p\) ，即仿射空间 \(\mathbf{A}_{\kappa(p)}^n = \text{Spec}(E(p))\) 中由 \(P(p)\) 所定义的子簇，而集合 \(F\) ——由使得该子簇为整个空间的 \(p\) （即使 \(P(p) = 0\) ）组成——是

Spec(A) 的闭子集。前述条件意味着该闭集为空，换言之，对每个 p，子簇 \(V_{p}\) 是 \(\mathbf{A}_{\kappa(p)}^{n}\) 中的一个超曲面。 *

\begin{enumerate}
\def\labelenumi{\arabic{enumi})}
\setcounter{enumi}{10}
\tightlist
\item
  设 S 和 X 为两个复解析空间，且 \(f: \mathbf{X} \to \mathbf{S}\) 为一个态射。我们说 \(f\) 在点 \(x\) （属于 \(\mathbf{X}\) ）处是平坦的，如果 \(\mathcal{O}_{\mathbf{X},x}\) 被视为 \(\mathcal{O}_{\mathbf{S},f(x)}\) -模（经由同态 \(f^{*}: \mathcal{O}_{\mathbf{S},f(x)} \to \mathcal{O}_{\mathbf{X},x}\) ）时是平坦的。\(\mathbf{X}\) 中使 \(f\) 平坦的点所成的集合是 \(\mathbf{X}\) 的一个开子集，并且 \(f\) 在这个开子集上的限制是一个开映射。如果 \(\mathbf{X}\) 和 \(\mathbf{S}\) 是连通的有限维解析流形，那么 \(f\) 是平坦的（在 \(\mathbf{X}\) 的每一点处）当且仅当 \(f(\mathbf{X})\) 在 \(\mathbf{S}\) 中是开的，并且纤维 \(f^{-1}(s)\) （对于 \(s \in f(\mathbf{X})\) ）都具有相同的维数。*
\end{enumerate}

